% -----------------------------------------------------------------------------------------------------
\section{全阶因子化 \texorpdfstring{$F=H\cdot J_1J_2/S$}{F = H J1 J2 / S}}\label{sec:sudakov-factorization}

一圈区域分析显示了全阶结构的雏形：硬区与红外无关，两个共线区各自只“认识”一条外线，超软区只“认识”外线的方向与虚度。全阶因子化要使这一图像对任意阶成立，并给出各因子的\emph{算符定义}。

\paragraph{领头区域与近似。}由\cref{thm:ir-leading-regions}（全阶论证，[A]），类时组态的领头约化图（\cref{fig:sudakov-hjjs}(a)）由硬子图 $H$、沿 $n_A$、$n_B$ 的喷注子图 $J_1$、$J_2$
与只以胶子连到喷注上的软子图 $S$ 组成；每个喷注以一个物理夸克加任意多个\emph{标量极化}共线胶子连到 $H$，没有软线连到 $H$。由\cref{prop:ir-power-counting}，这些构型 $\omega=0$
（多余的标量极化胶子不付出 $\lambda$ 的幂次），其余连接至少付出一个 $\lambda$（E-S 型陈述）。共线近似把 $J_1$ 与 $H$ 之间胶子（动量 $k$ 流入 $H$）的 $g_{\mu\nu}$ 换成
$K_{\mu\nu}(k)=k_\mu n_{B\nu}/(n_B\cdot k-\ii0)$（$\mu$ 在 $H$ 一侧），即 Collins 对 Grammer--Yennie 分解\cite{GrammerYennie1973}的单侧改写中的 K 部分\cite[式~(10.23)，第~10.6~节]{Collins2011}；
软近似在喷注一侧保留（投影后的）$k^\mu$、在软子图一侧换成 $n_A^\nu/(n_A\cdot k-\ii0)$（$k$ 由 $S$ 流入 $J_1$；$-\ii0$ 对应指向未来的 Wilson 线，\cite[式~(10.20)、(10.68)]{Collins2011}）。两者都把一个指标换成 $k^\mu$，于是可以用 Ward 恒等式。

\begin{example}[Ward 恒等式把纵向胶子收集为 Wilson 线]\label{ex:sudakov-ward-wilson}
取 $H$ 的最低阶：流顶点 $\Gamma$（$=\gamma^\rho$，通向 $J_1$ 的 Dirac 指标保持开放）与出射反夸克 $v(p_2)$；$J_1$ 中的 A 共线胶子（动量 $k$ 流入 $H$，色 $a$）附着在 $H$ 内的反夸克线上：
$H^\mu_a=\Gamma\,S(-p_2+k)\,(\ii g\gamma^\mu t^a)\,v(p_2)$，$S(p)=\ii\sla p/(p^2+\ii0)$。

(i) \emph{一个胶子。}由 $\sla p_2v(p_2)=0$，$\sla kv(p_2)=(\sla k-\sla p_2)v(p_2)$，故 $S(-p_2+k)\,\ii g\sla k\,v(p_2)=-g\,v(p_2)$（出射反夸克端的基本 Ward 恒等式），
对任意 $k$ 有 $k_\mu H^\mu_a=-g\,t^aH_0$，$H_0\equiv\Gamma\,v(p_2)$。% VF-WL-6 (W2)
$K_{\mu\nu}$ 中的 $-\ii0$ 取自被近似的传播子：$(p_2-k)^2+\ii0\simeq-2p_2^-(k^+-\ii0)$。于是 $H^\mu K_{\mu\nu}C^\nu=H_0\,[-g\,t^a\,n_B\!\cdot\!C/(n_B\cdot k-\ii0)]$，
方括号恰是 $W[0,\infty n_B]$ 的一阶项（\cref{prop:prelim-wilson-rules}：取向 $-n_B$，顶点 $-\ii gt^an_B^\nu$，传播子 $-\ii/(n_B\cdot k-\ii0)$）：反夸克线被一条沿 $n_B$、伸向未来无穷远的 Wilson 线取代。

(ii) \emph{$m$ 个胶子，阿贝尔情形。}对全部 $m!$ 种附着次序求和，内部插入由 $S(p+k)(\ii g\sla k)S(p)=g[S(p+k)-S(p)]$ 逐项相消，末次插入用 (i)，
得 $\sum_{\text{次序}}\prod_jk_{j\mu_j}H^{\mu_1\cdots\mu_m}=(-g)^mH_0$（对任意 $k_j$ 精确）；乘上 $m$ 个 K 投影，由程函恒等式（\cref{lem:prelim-eikonal-identity}）
$H_0\prod_j[-g\,n_B\cdot C_j/(n_B\cdot k_j-\ii0)]$ 正是阿贝尔 Wilson 线 $W[0,\infty n_B]$ 的 $m$ 阶项。

(iii) \emph{两个胶子，非阿贝尔树图。}对最后吸收的胶子用 (i)，再把 $\sla k_2t^bt^a+\sla k_1t^at^b$ 写成 $\sla Kt^at^b-\sla k_2[t^a,t^b]$（$K=k_1+k_2$），得精确恒等式
\[
k_{1\mu}k_{2\nu}H^{\mu\nu}_{ab}=g^2t^at^bH_0+\ii g^2[t^a,t^b]\,\Gamma S(-p_2+K)\sla k_2v(p_2);
\]
对严格共线的胶子 $k_j=k_j^+n_A$，第二项化为 $-g^2[t^a,t^b]\,(k_2^+/K^+)\,H_0$，两项合起来是 $(n_B\!\cdot\!k_1)(n_B\!\cdot\!k_2)\,W^{(2)}_{ab}[0,\infty n_B]\,H_0$：
对易子项恰好补出路径排序。DIS 中出射夸克的同一恒等式见附录~F 的\cref{prop:appF-two-gluon}；以上各式都以显式 Dirac 矩阵与 SU(3) 生成元核对过。% check: writing/work/W05/s6_ward_ff.py
\uses{def:prelim-feynman-rules,prop:prelim-wilson-rules,lem:prelim-eikonal-identity}
\end{example}

%%% BEGIN fig:sudakov-hjjs
\begin{figure}[htbp]
\centering
\begin{tikzpicture}[fd,
  fa lab/.style={font=\small},
  fa note/.style={font=\footnotesize},
  fa cap/.style={font=\footnotesize},
  fa eq/.style={font=\normalsize}]
  % ===================== (a) leading region (reduced diagram) =====================
  \begin{scope}
    \node[hard blob,circle,minimum size=1.0cm,inner sep=0pt] (H) at (0,0) {$H$};
    \node[jet blob,ellipse,minimum width=2.15cm,minimum height=0.66cm,rotate=28] (J1) at (2.8,1.55) {};
    \node[jetB blob,ellipse,minimum width=2.15cm,minimum height=0.66cm,rotate=-28] (J2) at (2.8,-1.55) {};
    \node[soft blob,circle,minimum size=0.62cm,inner sep=0pt] (S) at (3.4,0) {};
    \draw[photon] (-1.75,0) -- (H.west);
    \path[mom={$q$},mom sep=0.28cm] (-1.75,0) -- (H.west);
    % quark: H -> J1 -> out;  antiquark: out -> J2 -> H  (arrows = fermion-number flow)
    \draw[quark] (H.60) -- (J1.180);
    \draw[quark] (J1.0) -- ++(28:1.25) coordinate (E1);
    \draw[quark] ($(J2.0)+(-28:1.25)$) coordinate (E2) -- (J2.0);
    \draw[quark] (J2.180) -- (H.-60);
    % momenta of the outgoing quark and antiquark (both point away from H)
    \path[mom={$p_1$},mom pos=0.55] (J1.0) -- (E1);
    \path[mom'={$p_2$},mom pos=0.55] (J2.0) -- (E2);
    % scalar-polarized collinear gluons J_i -> H
    \draw[longit] (J1.246) to[bend left=8] (H.14);  \longitmark{(J1.246)}{(H.14)}
    \draw[longit] (J2.114) to[bend right=8] (H.-14); \longitmark{(J2.114)}{(H.-14)}
    % soft gluons, attached to the jets only
    \draw[gluon,green!45!black] (S.118) -- (J1.290);
    \draw[gluon,green!45!black] (S.62) -- (J1.320);
    \draw[gluon,green!45!black] (S.-118) -- (J2.70);
    \draw[gluon,green!45!black] (S.-62) -- (J2.40);
    \node[fa lab] at (2.55,2.27) {$J_1$};
    \node[fa lab] at (2.55,-2.27) {$J_2$};
    \node[fa lab,anchor=west] at (S.east) {$S$};
    \node[fa note,anchor=west] at ($(E1)+(0.05,0)$) {沿 $n_A$};
    \node[fa note,anchor=west] at ($(E2)+(0.05,0)$) {沿 $n_B$};
    \node[fa note,text=blue!55!black,align=center,anchor=east] at (0.55,-1.3) {纵向极化\\共线胶子};
    \draw[blue!55!black,line width=0.4pt] (0.5,-1.05) -- (0.98,-0.5);
    \node[fa cap] at (1.75,-3.0) {(a)\ 领头区域（约化图）};
  \end{scope}
  % ===================== (b) factorized form F = H J1 J2 / S =====================
  \begin{scope}[yshift=-5.1cm,xshift=-2.45cm]
    \node[fa eq] at (-0.75,0) {$F\;=\;H\;\times$};
    % ---- J1: quark jet along n_A, Wilson line along n_B ----
    \begin{scope}[xshift=0.45cm]
      \coordinate (O) at (0,0);
      \node[jet blob,ellipse,minimum width=1.3cm,minimum height=0.5cm,rotate=28] (J) at (28:1.4) {};
      \draw[quark] (O) -- (J.180);
      \draw[quark] (J.0) -- ++(28:0.7);
      \draw[wilson=0.22] (-28:2.1) coordinate (W) -- (O);
      \draw[gluon] (J.275) -- ($(O)!0.47!(W)$);
      \node[vertex] at ($(O)!0.47!(W)$) {};
      \node[opvertex] at (O) {};
      \node[fa note,anchor=west] at (W) {$n_B$};
      \node[fa lab] at (1.25,1.45) {$J_1$};
    \end{scope}
    \node[fa eq] at (3.15,0) {$\times$};
    % ---- J2: antiquark jet along n_B, Wilson line along n_A ----
    \begin{scope}[xshift=3.75cm]
      \coordinate (O) at (0,0);
      \node[jetB blob,ellipse,minimum width=1.3cm,minimum height=0.5cm,rotate=-28] (J) at (-28:1.4) {};
      \draw[quark] (J.180) -- (O);
      \draw[quark] ($(J.0)+(-28:0.7)$) -- (J.0);
      \draw[wilson=0.8] (O) -- (28:2.1) coordinate (W);
      \draw[gluon] (J.85) -- ($(O)!0.47!(W)$);
      \node[vertex] at ($(O)!0.47!(W)$) {};
      \node[opvertex] at (O) {};
      \node[fa note,anchor=west] at (W) {$n_A$};
      \node[fa lab] at (1.25,-1.45) {$J_2$};
    \end{scope}
    \node[fa eq] at (6.55,0) {$\Big/$};
    % ---- S: two Wilson lines meeting at a cusp ----
    \begin{scope}[xshift=7.05cm]
      \coordinate (O) at (0,0);
      \draw[wilson=0.84] (O) -- (28:2.1) coordinate (WA);
      \draw[wilson=0.16] (-28:2.1) coordinate (WB) -- (O);
      \node[soft blob,circle,minimum size=0.52cm,inner sep=0pt] (Sb) at (1.5,0) {};
      \draw[gluon,green!45!black] (Sb.110) -- ($(O)!0.5!(WA)$);
      \draw[gluon,green!45!black] (Sb.-110) -- ($(O)!0.5!(WB)$);
      \node[vertex] at ($(O)!0.5!(WA)$) {};
      \node[vertex] at ($(O)!0.5!(WB)$) {};
      \node[opvertex] at (O) {};
      \node[fa note,anchor=west] at (WA) {$n_A$};
      \node[fa note,anchor=west] at (WB) {$n_B$};
      \node[fa lab] at (1.25,1.45) {$S$};
    \end{scope}
    \node[fa cap] at (4.2,-1.85) {(b)\ 因子化形式 $F=H\,J_1J_2/S$};
  \end{scope}
\end{tikzpicture}
\caption{Sudakov 形状因子的领头区域与因子化 $F=H\,J_1J_2/S$（$\gamma^*\to q(p_1)\,\bar q(p_2)$，$p_1$ 沿 $n_A$、$p_2$ 沿 $n_B$）。
(a) 硬子图 $H$、两个喷注子图 $J_1$、$J_2$ 与只连接喷注的软子图 $S$（绿色）；Feynman 规范下 $J_i$ 与 $H$ 之间还有任意多条标量（纵向）极化的共线胶子（蓝色，以圆点标记）。
(b) 喷注函数中的 Wilson 线沿另一喷注方向，软函数是两条 Wilson 线的真空期望（$\otimes$ 为算符所在点，双线箭头与被替代的夸克线的费米子数流向一致）。}
\label{fig:sudakov-hjjs}
\end{figure}
%%% END fig:sudakov-hjjs

\cref{fig:sudakov-hjjs}(b) 画的是 Wilson 线取到光锥时的形式，$J_1$ 中的 $W[0,\infty n_B]$ 正是\cref{ex:sudakov-ward-wilson}得到的线；严格的陈述需要让 Wilson 线稍偏离光锥，以正规化\cref{ex:sudakov-massive}所示的快度发散。

\begin{theorem}[Sudakov 形状因子的全阶因子化 \tagA]\label{thm:sudakov-factorization}
考虑类时组态 $\gamma^*(q)\to q(p_1)\,\bar q(p_2)$ 的形状因子（\cref{def:sudakov-ff}），$p_1$ 沿 $n_A$、$p_2$ 沿 $n_B$ 运动，红外发散以规范不变的方式正规化
（阿贝尔规范理论中可取胶子质量与在壳有质量费米子，如\cite[第~10~章]{Collins2011}；QCD 中取维数正规化下的在壳无质量夸克）。取有限快度 $y_A>y_B$ 的类空辅助矢量
（光锥分量 $(v^+,v^-,\mathbf v_T)$）$\tilde n_A=(1,-\ee^{-2y_A},\mathbf 0_T)$、$\tilde n_B=(-\ee^{2y_B},1,\mathbf 0_T)$（$y_A\to+\infty$、$y_B\to-\infty$ 时趋于 $n_A$、$n_B$），
令 $P_B=\gamma^+\gamma^-/2$，$\gamma^\pm=(\gamma^0\pm\gamma^3)/\sqrt2$，并定义共线因子 $J_{1,2}$ 与软因子 $S$：
\begin{equation}\label{eq:sudakov-JS}
\begin{gathered}
\langle q(p_1)|\,\bar\psi(0)\,W[0,\infty\tilde n_B]\,P_B\,|0\rangle=J_1\,\bar u(p_1)P_B,\\
\langle\bar q(p_2)|\,W[\infty\tilde n_A,0]\,P_B\,\psi(0)\,|0\rangle=J_2\,P_Bv(p_2),\\
S=\frac1{N_c}\,\langle0|\,\Tr\,T\{W[\infty\tilde n_A,0]\,W[0,\infty\tilde n_B]\}\,|0\rangle ,
\end{gathered}
\end{equation}
其中色结构为单位矩阵，各矩阵元都已除以相应 Wilson 线的自能因子并作 $\MSbar$ 紫外重整化（\cite[式~(10.89)--(10.92)、(10.101)]{Collins2011}）。则在领头幂次、对 $\alpha_s$ 的所有阶，
\begin{equation}\label{eq:sudakov-fact}
F=H\Big(\frac{Q^2}{\mu^2},\frac{\zeta_1}{\mu^2},\frac{\zeta_2}{\mu^2};\alpha_s(\mu)\Big)\,\frac{J_1\,J_2}{S}+(\text{幂次压低项}),\qquad
\zeta_1=2(p_1^+)^2\ee^{-2y_B},\quad\zeta_2=2(p_2^-)^2\ee^{2y_A},
\end{equation}
其中 $J_i$ 只依赖于 $p_i$、相应的辅助矢量与红外正规子，$S$ 只依赖于 $y_A-y_B$ 与红外正规子，硬因子 $H$ 与红外正规子无关。
等价地（Collins 的写法），$F=\lim H\,J_1(y_{p_1}-y_{u})\,J_2(y_{u'}-y_{p_2})\,S(y_A-y_B)/[S(y_{u'}-y_B)\,S(y_A-y_u)]$，极限取 $y_{u'}\to+\infty$、$y_u\to-\infty$，
快度发散在分子分母之间相消\cite[式~(10.109)]{Collins2011}；两种写法中的 $H$ 只相差一个仅含硬动量的因子。（本定理只涉及部分子态。）
% errtype: E-S
\uses{def:sudakov-ff,thm:ir-leading-regions,prop:ir-power-counting,lem:prelim-eikonal-identity,prop:prelim-wilson-rules,def:prelim-wilson,thm:prelim-slavnov-taylor,prop:sudakov-oneloop,ex:sudakov-offshell,ex:sudakov-ward-wilson}
\end{theorem}

\begin{proofsketch}
(1) \emph{领头区域}：由\cref{thm:ir-leading-regions}与\cref{prop:ir-power-counting}（E-S 型陈述）。
(2) \emph{近似}：共线近似用类光矢量，$k_\mu n_{B\nu}/(n_B\cdot k-\ii0)$；软近似在软子图一侧用有限快度的 $\tilde n_A^\nu/(\tilde n_A\cdot k-\ii0)$，以正规化快度发散\cite[式~(10.15)--(10.24)]{Collins2011}。
各近似在其区域的核心幂次精确，但当内禀变量趋于更小区域、或（对软胶子）角变量趋于更大区域的钳夹面时精度下降，这正是减除的原因\cite[第~10.3、10.4、10.6~节]{Collins2011}。
(3) \emph{Ward 恒等式}：阿贝尔情形见\cref{ex:sudakov-ward-wilson}（工具为\cref{prop:prelim-wilson-rules}、\cref{lem:prelim-eikonal-identity}），纵向共线胶子给出 $J_i$ 中的 Wilson 线；
同一论证作用于喷注一侧，把软胶子从喷注上解耦并连到 $S$ 的 Wilson 线上。Collins 在协变规范的阿贝尔理论中连同减除项与不可约性要求对所有阶完成了这一步\cite[第~10.8.3--10.8.5~节]{Collins2011}；
非阿贝尔情形还需\cref{thm:prelim-slavnov-taylor}与含鬼场的组合学（树图见\cref{ex:sudakov-ward-wilson}(iii)，全阶形式见第~9~章\cref{lem:method-ward-multi}，前向指针）。
(4) \emph{减除}：一圈时 $F=H+C_1+C_2+S$ 且重叠无标度（\cref{ex:sudakov-offshell}），正是 $H\cdot J_1J_2/S$ 的一圈展开。全阶上，Collins 的逐次近似
$C_R=T_R(1-\sum_{R'<R}C_{R'})$\cite[式~(10.4)]{Collins2011}对区域与图求和后，减除项以“除以软因子”的形式出现\cite[式~(10.94)--(10.95)]{Collins2011}，给出 $F=H\,ABS$\cite[式~(10.11)]{Collins2011}，
即上述极限形式\cite[式~(10.109)]{Collins2011}。快度演化方程\cite[式~(10.112)--(10.117)]{Collins2011}表明 $J_1(y_{p_1}-y_u)/S(y_A-y_u)$ 对 $y_u$ 的依赖只来自 $\frac12G$（$K$ 项在比值中相消），
它只含硬动量而可吸收进 $H$，故可取 $y_u=y_B$、$y_{u'}=y_A$，得式~\eqref{eq:sudakov-fact}（$H$ 因而依赖于 $\zeta_1,\zeta_2$）。
一般的减除定理与此三角形的单圈减除见第~9~章\cref{thm:method-subtraction}与\cref{ex:method-sudakov-subtraction}（前向指针）。
完整论证见\cite[第~10.3、10.8--10.10~节]{Collins2011}（阿贝尔情形），QCD 中见\cite{Mueller:1979ih,Collins:1980ih,Sen:1981sd,Korchemsky1989,Collins:1989bt,Magnea:1990zb}，宽角振幅的相应分析见\cite{Akhoury1979,Sen1983}。
\end{proofsketch}

\paragraph{地位与缺口}本定理是全阶论证（[A]），其瓶颈有两处。其一是\emph{减除}：Collins~2011 第~10~章对此例给出的减除论证（误差估计即\cite[第~10.7~节，式~(10.81)--(10.84)]{Collins2011}的局部幂次压低定理）在该章的严格程度下是完整的，
但它停留在标度陈述与钳夹面的层次（E-S），而没有给出积分范数界（E-N）；积分范数下的估计只对单圈离壳三角形写出（附录~N 的\cref{prop:appN-triangle-cs}）。
因此 $[A]$ 等级的瓶颈是\cref{asm:appN-collins}（全阶论证，[A]）的内容，其精确的条件形式是\cref{thm:appN-subtraction}——这里只作指针而不作前提，因为两者都是用第~9~章的条件\cref{def:method-collins-conditions}表述的。
其二是\emph{非阿贝尔 Ward 恒等式}：Collins 的全阶推导限于阿贝尔理论（\cite[第~10.8.3~节]{Collins2011}），QCD 情形依赖第~9~章\cref{lem:method-ward-multi}（[A]）的多胶子 Ward 恒等式，
本书在那里点名该论证中被明确假定的一步（\cref{asm:appN-brst-nonlocal}，[C]，具名缺口）。领头区域一步继承\cref{thm:ir-leading-regions}的 E-S 地位。
把本定理升级为 $[R]$ 需要：一般图在积分范数下的减除估计（\cref{asm:appN-collins}），以及含鬼场的非阿贝尔多胶子 Ward 恒等式的完整组合学证明。
